\documentclass[11pt,a4paper]{article}
\usepackage[margin=2.2cm]{geometry}
\usepackage{amsmath,amssymb,booktabs,longtable,xcolor,hyperref}
\usepackage[T1]{fontenc}
\hypersetup{colorlinks=true,linkcolor=blue,urlcolor=blue}
\newcommand{\code}[1]{\texttt{\small #1}}
\setlength{\parskip}{0.35em}

\title{\vspace{-1.5em}\textbf{It was never the geometry:}\\[0.2em]
       \large the remapped \emph{dynamics seam} is what kills FESOM on a moving-cavity leg\\[0.3em]
       \normalsize\textit{Findings of 2026-07-13 --- supersedes the balance-consistent restart plan}}
\author{}
\date{2026-07-13}

\begin{document}
\maketitle
\vspace{-2em}

\begin{abstract}
\noindent
This note reports a set of controlled experiments that \textbf{falsify the previous root-cause
claim} of this investigation. The moving-cavity blowup is \emph{not} a geostrophically
unbalanced restart, \emph{not} the basal-melt flux shock, \emph{not} the timestep, and
\emph{not} the cavity geometry. A \textbf{cold start on the very same PISM cavity, at the full
production timestep, is stable}; the identical mesh warm-started from our remapped restart dies
on day~4. What kills the model is the \textbf{seam between the evolved dynamics we carry over
on the untouched 98.2\% of the mesh and the patched values at the changed nodes}. We also show
that the change we have been testing all along --- the one-off CORE3(observed)$\to$PISM cavity
swap --- is a pathological monster ($\max|\Delta\mathrm{draft}| = 2711$~m) that production never
has to survive, whereas the change production \emph{does} have to survive, a 10-yr PISM
increment, has \textbf{never been tested}. The strategy therefore changes: cold-start the ocean
on PISM's cavity and never remap across the swap at all.
\end{abstract}

\section{What is now falsified (including our own previous conclusion)}

The design note \code{balanced\_restart\_plan.tex} argued the blowup was a missing geostrophic
jet and proposed initialising the changed-cell velocity in thermal-wind balance. That was
implemented (per-level $\mathbf{u}_g = (1/f\rho_0)\,\hat{z}\times\nabla p$ using FESOM's own
cavity pressure convention). It \textbf{failed its own offline gate}: the discrete $\nabla p$ at
grid-scale fronts is \emph{noisier} than the nearest-neighbour fill (transport divergence above
the NN baseline), so it ships \textbf{opt-in only} (\code{REMAP\_GEOSTROPHIC=1}) and is off by
default. It is not the fix, and the hypothesis behind it is dead.

Every hypothesis we have held is now eliminated by a direct experiment:

\begin{longtable}{@{}p{5.0cm}p{6.6cm}p{3.4cm}@{}}
\toprule
\textbf{Hypothesis} & \textbf{Experiment} & \textbf{Verdict} \\
\midrule
\endhead
Atmosphere / OASIS / coupling & ocean-only standalone FESOM reproduces the crash exactly
  (day~4, CFLz~9.24) & \textbf{ruled out} \\
CORE2 forcing shock (standalone artefact) & full-mesh control, \emph{same} forcing + native
  restart: stable 21~d, worst CFLz \textbf{2.61} & \textbf{ruled out} \\
Remap corrupting the state & remapped vs.\ source over the 204\,970 unchanged nodes:
  $\max|\Delta| = \mathbf{0.000}$ for \code{temp}, \code{salt}, \code{ssh}, \code{hnode} &
  \textbf{ruled out} (bit-exact) \\
The restart \emph{fill values} at changed nodes & five independent fill strategies
  (v2--v8: donor, coherent-column, shelf-base hold, NN velocity, full balance hygiene) &
  \textbf{ruled out} (all die $\approx$day 4) \\
Pathological thin water columns & submesh min wet column $30$~m vs.\ mother's $25$~m;
  min 3 wet layers on changed nodes; zero 1--2-layer columns & \textbf{ruled out}
  (submesh is \emph{cleaner}) \\
Basal melt flux shock & \code{cavity\_flux\_ramp\_days} implemented; melt effectively
  \textbf{OFF} still dies (mstep 328) & \textbf{ruled out} \\
Timestep / unresolved $w$ & 1200\,s: CFLz kill day 4.0. 60\,s: \textbf{0} CFLz warnings, dies
  day 4.4 via $\eta_n$. 120\,s: dies day 8.9 & \textbf{ruled out} \\
Sub-ice free-surface BC violation & real bug found + fixed (987 nodes, below); run still
  dies day 4.4 & \textbf{real bug, not the cure} \\
\textbf{Cavity geometry itself} & \textbf{cold start on the same PISM submesh @1200\,s:
  STABLE, day 13+, worst CFLz 3.44} & \textbf{ruled out} \\
\bottomrule
\end{longtable}

\section{The decisive control: a cold start on PISM's cavity is stable}\label{sec:cold}

The single experiment that inverts the picture. Same submesh, same forcing, same production
timestep --- only the initial state differs:

\begin{center}
\begin{tabular}{@{}llll@{}}
\toprule
Run & $\Delta t$ & Initial state & Outcome \\
\midrule
full-mesh control & 1200\,s & native spun-up CORE3 restart & stable 21\,d, worst CFLz \textbf{2.61} \\
\textbf{PISM cavity, cold} & \textbf{1200\,s} & \textbf{cold (climatology, at rest)} &
  \textbf{stable 13+\,d, worst CFLz 3.44} \\
PISM cavity, warm & 1200\,s & our remapped restart & \textbf{dies day 4.0}, CFLz 9.24 \\
PISM cavity, warm & \phantom{0}120\,s & our remapped restart & dies day 8.9 ($\eta_n$) \\
PISM cavity, warm & \phantom{00}60\,s & our remapped restart & dies day 4.4 ($\eta_n$, 0 CFLz) \\
\bottomrule
\end{tabular}
\end{center}

\noindent
\textbf{The ocean has no difficulty living in PISM's cavity at the production timestep.} Its
CFLz (3.44) sits essentially at the full-mesh control's background level (2.61). Whatever is
wrong is in \emph{our restart}, not in the mesh, the physics, or the coupling.

\section{What is actually wrong: the dynamics seam}\label{sec:seam}

The remap produces a state that is internally \emph{inconsistent}:

\begin{itemize}
\item on the \textbf{204\,970 unchanged nodes} (98.2\%) it reproduces the evolved restart
      \textbf{bit-exactly} --- a fully developed, geostrophically balanced dynamic state
      ($\mathbf{u}$, $\eta$, AB history);
\item on the \textbf{3840 changed nodes} it substitutes \emph{patched} values (donor-filled
      tracers, NN or zero velocity, donor $\eta$).
\end{itemize}

\noindent
Those two regions are mutually incompatible, and the \textbf{seam between them} is the
instability. The evidence is consistent and specific:

\begin{itemize}
\item the CFLz runaway grows \emph{monotonically} ($4.46 \to 6.43 \to 7.25 \to 7.51 \to 7.92
      \to 9.24$);
\item every escalation site lies \textbf{within $0.16$--$0.22^\circ$ of a changed node}, and
      the worst one \emph{is} a changed cavity node --- it grows in the ring around the patch
      and spreads;
\item it is independent of melt (survives melt-OFF) and of $\Delta t$ (60\,s merely swaps the
      \emph{symptom} from CFLz to $\eta_n$);
\item \textbf{zeroing the dynamics globally} (keep remapped $T,S$; set
      $\mathbf{u}=\eta=\bar{h}=w=\mathrm{AB}=0$, \code{hnode} nominal) \textbf{removes the
      seam and the crash} --- the ocean then re-derives its own balanced flow from rest, exactly
      as the cold start does.
\end{itemize}

\noindent
This finally explains why \emph{five} fill strategies all failed identically: each one varied
the \emph{patch} while leaving the \emph{seam} in place. The problem was never the values we
put at the changed nodes; it is that we splice \emph{any} patch against a fully evolved
dynamical field.

\medskip\noindent
\textbf{But zero-dynamics is not a production fix.} Discarding the ocean's circulation at every
coupling leg is unacceptable --- the whole point is to remap restarts for an ever-evolving PISM
geometry and have FESOM \emph{survive with its dynamics intact}. It is a diagnostic, not a cure.

\section{A real bug found along the way (not the cure)}\label{sec:ssh}

FESOM keeps $\eta \equiv 0$ under an ice shelf: in the source restart, every already-sub-ice
node is \emph{exactly} $0.000$. The remap, however, carried the old \emph{open-ocean} value
across for nodes the mesh change puts \textbf{newly under ice}:

\begin{center}
\begin{tabular}{@{}lrr@{}}
\toprule
node class & count & \code{ssh} \\
\midrule
still-ice (ice before \emph{and} after) & 2368 & $+0.0000$ (exactly) \\
\textbf{newly-ICE} (open before, sub-ice now) & \textbf{987} & $\mathbf{-1.60}$~m (all nonzero) \\
open-both & 204\,797 & $-0.46$ (normal range) \\
\bottomrule
\end{tabular}
\end{center}

\noindent
987 nodes entered the ssh solver with a free surface it can never satisfy --- a standing
violation of the model's own surface BC, with \code{hbar} corrupted identically. This is a
genuine correctness bug and is \textbf{fixed} (zero \code{ssh}/\code{hbar} at
\code{ulevels\_nod2D > 1}). It is \textbf{not} the cure: the fixed run still dies on day 4.4.
We report it because it was a plausible-looking culprit that explained much of the evidence
(melt-independent, $\Delta t$-independent, barotropic, margin-localised) and \emph{was still
wrong} --- a caution for the next hypothesis.

\section{We were testing the wrong change}\label{sec:wrongchange}

The most consequential realisation. The mesh change we have spent this entire investigation
fighting is the \textbf{one-off CORE3(observed)$\,\to\,$PISM cavity swap}, which is
pathological:

\begin{center}
\begin{tabular}{@{}lr@{}}
\toprule
& magnitude \\
\midrule
$\max |\Delta \mathrm{draft}|$ (observed $\to$ PISM) & \textbf{2711\,m} \\
nodes with $|\Delta \mathrm{draft}| > 200$\,m & 1853 \\
nodes flipping open-ocean $\to$ sub-ice & 1579 \\
nodes flipping sub-ice $\to$ open-ocean & 1436 \\
\midrule
\textbf{a genuine 10-yr PISM increment}: mean $|\Delta \mathrm{draft}|$ & \textbf{$\approx$22\,m} \\
\bottomrule
\end{tabular}
\end{center}

\noindent
This swap is an artefact of bolting an ocean spun up on the \emph{observed} cavity onto PISM's
\emph{equilibrium} cavity --- two different ice sheets. \textbf{Production never has to survive
it.} What production must survive, forever, is the \emph{incremental} change from one PISM
chunk to the next --- two orders of magnitude smaller --- and \textbf{that has never been
tested}.

\section{The strategy that follows}\label{sec:strategy}

\begin{enumerate}
\item \textbf{Never remap across the swap.} \textbf{Cold-start} FESOM \emph{on PISM's cavity}
      in chunk~1 (\S\ref{sec:cold} proves this is stable at 1200\,s). The observed cavity never
      enters the coupled run. This is experiment \textbf{movcav12} (running).
\item \textbf{Then test the change that matters.} By chunk~3 the cold-started ocean carries a
      year of genuine, spun-up dynamics; the remap must then carry \emph{that} across a real
      10-yr PISM increment ($\approx$22\,m). \textbf{This is the experiment that decides whether
      the coupling is viable}, and it is the first honest test of the remap we will have run.
\item \textbf{If the increment survives}, the machinery is fit for purpose --- we have already
      proven the remap is bit-exact on unchanged nodes, inversion-free, and geometrically sound.
\item \textbf{If even $\approx$22\,m dies}, the seam is fundamental and the remaining option is
      to stage the geometry change so the ocean never sees a discontinuity (a geometry ramp), or
      to relax the dynamics locally toward balance across the seam rather than patching it.
\end{enumerate}

\section{Assets built (reusable)}

\begin{itemize}
\item \textbf{A 2-minute offline reproducer}: standalone ocean-only FESOM (\code{setup\_name:
      fesom} via esm\_tools) on the PISM-cavity submesh, 14 nodes, crashes in $\sim$2\,min ---
      versus the 45-node coupled system. Runscript:
      \code{scripts\_is/sa\_spin\_movcav11.yaml}. Requires a \emph{standalone} binary
      (\code{build\_sa}, \code{FESOM\_COUPLED=OFF}) rebuilt at the current SHA; note FESOM puts
      its code in \code{libfesom.so} and \code{fesom.x} is a thin wrapper whose timestamp never
      changes on rebuild.
\item \textbf{A stable control}: full mesh + native restart + same forcing (21\,d, CFLz 2.61) ---
      the reference every experiment is measured against.
\item \code{cavity\_flux\_ramp\_days} (FESOM \code{\&run\_config}, default \code{0.0} = off,
      bit-identical): ramps cavity basal heat/freshwater fluxes $0\to1$ over $N$ model days.
      Melt was exonerated, but the knob is a clean, off-by-default spin-up device.
\item Sub-ice \code{ssh}/\code{hbar} zeroing in the remap (\S\ref{sec:ssh}).
\item \code{verify\_balance.py}: offline balance/stability gate (min $N^2$,
      $\max|\nabla\!\cdot\!(H\bar{\mathbf{u}})|$, $\max|\mathbf{u}|$, CFL, Ri) --- this is what
      caught the geostrophic $\nabla p$ init being worse than the NN baseline \emph{before} it
      cost a model run.
\end{itemize}

\section{Method note}

Two conclusions in this investigation were asserted with more confidence than the evidence
supported (``geostrophic imbalance is the root cause''; ``the sub-ice ssh BC violation is the
root cause''). Both were later falsified by a direct control. The controls that actually cracked
it --- the \textbf{full-mesh control} (is it our setup or the cavity?) and the \textbf{cold
start} (is it the geometry or our state?) --- are both cheap, and both should have been run
\emph{before} building fixes on top of an unverified hypothesis.

\end{document}
